\documentclass[12pt]{article}

\usepackage{amsmath}
\usepackage{amssymb}
\usepackage{physics}
\usepackage{hyperref}

\newcommand{\ds}{\displaystyle}

\title{Fields of a Uniformly Moving Charge}
\author{}
\date{}

\begin{document}

\maketitle


% ============================================================
\section{Introduction}
% ============================================================
We derive the electromagnetic field generated by a point charge moving with
constant velocity. Rather than starting from the known Lienard--Wiechert
potentials, we derive the result directly from Maxwell's equations by
introducing the scalar and vector potentials and solving the resulting wave
equations.

Throughout this section we use Gaussian units.

We begin with Maxwell's equations in vacuum,

\begin{align}
  \div \vb{E}
    &= 4\pi \rho, \\
  \div \vb{B}
    &= 0, \\
  \curl \vb{E}
    &= -\frac{1}{c}\pdv{\vb{B}}{t}, \\
  \curl \vb{B}
    &= \frac{1}{c}\left(4\pi\vb{J}+\pdv{\vb{E}}{t}\right).
    \label{EQ:MAXWELL}
\end{align}

These equations must be solved for the charge and current densities
corresponding to a point charge moving with constant velocity $\vb{v}$.
Assuming that the charge passes through the origin at $t=0$, the source terms
are

\begin{align}
  \rho(\vb{r},t)
    &= q\,\delta^3(\vb{r}-\vb{v}t), \\
  \vb{J}(\vb{r},t)
    &= q\,\vb{v}\,\delta^3(\vb{r}-\vb{v}t).
    \label{eq:uniform_motion_sources}
\end{align}

A convenient way to solve Maxwell's equations is to introduce a scalar
potential $\phi$ and a vector potential $\vb A$ defined by

\begin{align}
  \vb{E}
    &= -\nabla \phi - \frac{1}{c}\pdv{\vb{A}}{t}, \\
  \vb{B}
    &= \curl \vb{A}.
\end{align}

At this stage $\phi$ and $\vb A$ are unknown functions that must be determined
by solving Maxwell's equations. Their introduction is motivated by the fact
that Maxwell's equations acquire a much simpler form when expressed in terms
of the electromagnetic potentials.

The second Maxwell equation is automatically satisfied since

\[
  \div \vb{B}
  =
  \div(\curl\vb{A})
  =
  0.
\]

The third Maxwell equation is also identically satisfied. Indeed,

\begin{align}
  \curl \vb{E}
    &=
    \curl\left(
      -\nabla\phi
      -\frac{1}{c}\pdv{\vb{A}}{t}
    \right) \\
    &=
    -\curl(\nabla\phi)
    -\frac{1}{c}\pdv{}{t}(\curl\vb{A}) \\
    &=
    -\frac{1}{c}\pdv{\vb{B}}{t},
\end{align}

since $\curl\nabla\phi=\vb 0$.

Hence, after introducing the scalar and vector potentials, the two homogeneous
Maxwell equations are automatically satisfied.

We now substitute the potentials into Ampere's law. Using the identity

\[
\nabla\times(\nabla\times\mathbf A)
=
\nabla(\nabla\cdot\mathbf A)
-
\nabla^2\mathbf A,
\]

we obtain

\[
\nabla(\nabla\cdot\mathbf A)
-
\nabla^2\mathbf A
=
\frac{1}{c}
\left(
4\pi\mathbf J
+
\frac{\partial\mathbf E}{\partial t}
\right).
\]

Since

\[
\mathbf E
=
-\nabla\phi
-
\frac{1}{c}
\frac{\partial\mathbf A}{\partial t},
\]

its time derivative is

\[
\frac{\partial\mathbf E}{\partial t}
=
-
\nabla\frac{\partial\phi}{\partial t}
-
\frac{1}{c}
\frac{\partial^2\mathbf A}{\partial t^2}.
\]

Substituting this expression into the previous equation yields

\[
\nabla(\nabla\cdot\mathbf A)
+
\frac{1}{c}
\nabla\frac{\partial\phi}{\partial t}
=
\nabla^2\mathbf A
-
\frac{1}{c^2}
\frac{\partial^2\mathbf A}{\partial t^2}
+
\frac{4\pi}{c}\mathbf J.
\]

Our goal is to obtain an uncoupled wave equation for the vector potential.
This is achieved by requiring $\vb A$ to satisfy

\begin{equation}
\nabla^2\mathbf A
-
\frac{1}{c^2}
\frac{\partial^2\mathbf A}{\partial t^2}
=
-
\frac{4\pi}{c}\mathbf J,
\label{eq:wave_vector_potential}
\end{equation}

which implies that

\[
\nabla
\left(
\nabla\cdot\mathbf A
+
\frac{1}{c}
\frac{\partial\phi}{\partial t}
\right)
=
\mathbf0.
\]

The quantity inside the parentheses is therefore independent of position.
Since an arbitrary constant may be absorbed into the definition of the
potentials, we choose it to vanish and obtain the Lorenz gauge,

\begin{equation}
\nabla\cdot\mathbf A
+
\frac{1}{c}
\frac{\partial\phi}{\partial t}
=
0.
\label{eq:lorenz_gauge}
\end{equation}

The wave equation satisfied by the scalar potential follows from Gauss's law,

\[
-\nabla^2\phi
-
\frac{\partial}{\partial t}
\left(
\nabla\cdot\mathbf A
\right)
=
4\pi\rho.
\]

Substituting the Lorenz gauge condition,

\[
\nabla\cdot\mathbf A
=
-
\frac{1}{c}
\frac{\partial\phi}{\partial t},
\]

gives

\[
-\nabla^2\phi
+
\frac{1}{c^2}
\frac{\partial^2\phi}{\partial t^2}
=
4\pi\rho,
\]

or equivalently,

\begin{equation}
\nabla^2\phi
-
\frac{1}{c^2}
\frac{\partial^2\phi}{\partial t^2}
=
-4\pi\rho.
\label{eq:wave_scalar_potential}
\end{equation}

We have therefore transformed Maxwell's equations into two wave
equations for the scalar and vector potentials,

\begin{equation}
\left\{
\begin{aligned}
\nabla^2\mathbf A
-
\frac{1}{c^2}
\frac{\partial^2\mathbf A}{\partial t^2}
&=
-\frac{4\pi}{c}\mathbf J,
\\[0.6em]
\nabla^2\phi
-
\frac{1}{c^2}
\frac{\partial^2\phi}{\partial t^2}
&=
-4\pi\rho.
\end{aligned}
\right.
\label{eq:potential_wave_equations}
\end{equation}

supplemented by the Lorenz gauge
condition \eqref{eq:lorenz_gauge}, which couples the scalar and vector
potentials.

% ============================================================
\section{Green Function Solution}
% ============================================================

The two equations
\begin{align}
  \nabla^2\phi
  -
  \frac{1}{c^2}\frac{\partial^2 \phi}{\partial t^2}
  &=
  -4\pi\rho,
  \\
  \nabla^2\vb{A}
  -
  \frac{1}{c^2}\frac{\partial^2 \vb{A}}{\partial t^2}
  &=
  -\frac{4\pi}{c}\vb{J}
\end{align}
can be solved using the retarded Green function. In the present case, however,
it is sufficient to solve only the equation for the scalar potential. Indeed,
from \eqref{eq:uniform_motion_sources} the current density is simply
\begin{equation}
  \vb J=\rho\,\vb v,
\end{equation}
where the velocity $\vb v$ is constant. Consequently, if $\phi$ is the solution
of the scalar equation, the vector potential is immediately obtained as
\begin{equation}
  \vb A
  =
  \frac{\vb v}{c}\,\phi.
\end{equation}

Hence we proceed to find the solution for the scalar potential from the retarded Green function
\begin{equation}
  \phi(\vb{r},t)
  =
  \int
  \frac{
    \rho\!\left(
      \vb{r}',
      t-\frac{\abs{\vb{r}-\vb{r}'}}{c}
    \right)
  }{
    \abs{\vb{r}-\vb{r}'}
  }
  \,d^3 r'.
  \label{eq:scalar_potential_green}
\end{equation}

Here $\vb r'$ is the integration variable.  

For the uniformly moving point charge,
\[
  \rho(\vb{r}',t')
  =
  q\,\delta^3(\vb{r}'-\vb{v}t'),
\]
and therefore
\begin{equation}
  \phi(\vb{r},t)
  =
  q
  \int
  \frac{
    \delta^3\!\left[
      \vb{r}'
      -
      \vb{v}
      \left(
        t-\frac{\abs{\vb{r}-\vb{r}'}}{c}
      \right)
    \right]
  }{
    \abs{\vb{r}-\vb{r}'}
  }
  \,d^3 r'.
  \label{eq:scalar_potential_delta}
\end{equation}

The task is now to evaluate this integral. The difficulty is that the argument
of the delta function contains the integration variable $\vb{r}'$ not only explicitly, but also through
the distance $\abs{\vb{r}-\vb{r}'}$.

Therefore, in order to find the region of space where the argument of the $\delta$ vanishes, we need to solve the equation
\begin{equation}
  \vb{r}'
  -
  \vb{v}
  \left(
    t-\frac{\abs{\vb{r}-\vb{r}'}}{c}
  \right)
  =
  \vb{0}.
  \label{eq:source_position_condition}
\end{equation}

Given the particular form of this equation, it is useful to decompose the position vector into components parallel and
perpendicular to the velocity. Let us define the unitary vector in the direction of the particle velocity, 
\[
  \vb{n}
  =
  \frac{\vb{v}}{\abs{\vb{v}}},
\]

and decompose the position vectors in component parallel and perpendicolar to $\vb n$,

\[
  \vb{r}
  =
  r_\parallel \vb{n}
  +
  r_\perp \vb{n}_\perp,
  \qquad
  \ip{\vb{n}_\perp}{\vb{n}}=0.
\]

Similarly,
\[
  \vb{r}'
  =
  r'_\parallel \vb{n}
  +
  r'_\perp\vb{n}_\perp.
\]

From \eqref{eq:source_position_condition} we immediately conclude that
$\vb r'$ can only have a component parallel to the velocity. Therefore
\[
  \vb r'_\perp=\vb0.
\]

The vector equation \eqref{eq:source_position_condition} then reduces to the
scalar equation
\begin{equation}
  r'_\parallel
  -
  v
  \left(
    t
    -
    \frac{
      \sqrt{(r_\parallel-r'_\parallel)^2+r_\perp^2}
    }{c}
  \right)
  =
  0,
  \label{eq:source_position_parallel}
\end{equation}
 
We define
\[
  s=r'_\parallel-r_\parallel,
  \qquad
  \mu=r_\parallel-vt.
\]

Then equation \eqref{eq:source_position_parallel} becomes
\[
  s+\mu
  +
  v
  \frac{
    \sqrt{s^2+r_\perp^2}
  }{c}
  =
  0,
\]
or equivalently,
\[
  s+\mu
  =
  -
  \frac{v}{c}
  \sqrt{s^2+r_\perp^2}.
\]

Since the square root is non-negative, every solution of this equation must satisfy
\[
  s+\mu\le0.
\]
Under this condition the equation may be squared, yielding
\[
  s^2+\mu^2+2\mu s
  =
  \frac{v^2}{c^2}
  \left(
    s^2+r_\perp^2
  \right).
\]

Thus
\[
  \left(
    1-\frac{v^2}{c^2}
  \right)s^2
  +
  2\mu s
  +
  \mu^2
  -
  \frac{v^2}{c^2}r_\perp^2
  =
  0.
\]

This equation has the two algebraic solutions
\begin{equation}
  s
  =
  \frac{
    -\mu
    \pm
    \dfrac{v}{c}
    \sqrt{
      \left(
        1-\frac{v^2}{c^2}
      \right)r_\perp^2
      +
      \mu^2
    }
  }{
    1-\dfrac{v^2}{c^2}
  }.
  \label{eq:s_algebraic_solutions}
\end{equation}

However, not both algebraic solutions are admissible. 

We now prove which of the two algebraic solutions is admissible.  
Let
\[
  \alpha
  =
  \sqrt{
    \left(
      1-\frac{v^2}{c^2}
    \right)r_\perp^2
    +
    \mu^2
  }.
\]
Then the two roots can be written as
\[
  s
  =
  \frac{
    -\mu
    \pm
    \dfrac{v}{c}\alpha
  }{
    1-\dfrac{v^2}{c^2}
  }.
\]

Adding $\mu$ gives
\begin{align}
  s+\mu
  &=
  \frac{
    -\mu
    \pm
    \dfrac{v}{c}\alpha
  }{
    1-\dfrac{v^2}{c^2}
  }
  +
  \mu
  \\
  &=
  \frac{
    \pm\dfrac{v}{c}\alpha
    -
    \mu\dfrac{v^2}{c^2}
  }{
    1-\dfrac{v^2}{c^2}
  }.
\end{align}

Since $1-v^2/c^2>0$, the admissibility condition is equivalent to
\[
  \pm\frac{v}{c}\alpha
  -
  \mu\frac{v^2}{c^2}
  <
  0.
\]
Equivalently,
\[
  \pm\alpha
  -
  \mu\frac{v}{c}
  <
  0.
\]

For the root with the plus sign this would require
\[
  \alpha
  <
  \mu\frac{v}{c}.
\]
If $\mu<0$, this is impossible because $\alpha\ge0$ while
$\mu v/c<0$. If $\mu>0$, squaring would imply
\[
  \alpha^2
  <
  \mu^2\frac{v^2}{c^2}.
\]
But
\[
  \alpha^2
  =
  \left(
    1-\frac{v^2}{c^2}
  \right)r_\perp^2
  +
  \mu^2,
\]
so this would require
\[
  \left(
    1-\frac{v^2}{c^2}
  \right)r_\perp^2
  <
  \mu^2
  \left(
    \frac{v^2}{c^2}-1
  \right),
\]
which is impossible, because the left-hand side is non-negative while the
right-hand side is negative. Therefore the root with the plus sign is not
admissible.

For the root with the minus sign, the condition becomes
\[
  -\alpha
  -
  \mu\frac{v}{c}
  <
  0.
\]
If $\mu>0$, this is immediately satisfied. If $\mu<0$, it is
equivalent to
\[
  \alpha
  >
  -\mu\frac{v}{c},
\]
which follows from
\[
  \alpha^2
  =
  \left(
    1-\frac{v^2}{c^2}
  \right)r_\perp^2
  +
  \mu^2
  >
  \mu^2\frac{v^2}{c^2}.
\]

Thus the only admissible solution is the root with the minus sign, 
\begin{equation}
  s
  =
  \frac{
    -\mu
    -
    \dfrac{v}{c}
    \sqrt{
      \left(
        1-\frac{v^2}{c^2}
      \right)r_\perp^2
      +
      \mu^2
    }
  }{
    1-\dfrac{v^2}{c^2}
  }.
  \label{eq:s_admissible_solution}
\end{equation}

 Introducing the standard notation
\[
  \beta=\frac{v}{c},
\]
and expressing the result in terms of the original variables, we obtain
\begin{equation}
  r'_\parallel
  =
  \frac{
    \left(1-\beta^2\right)\mu
    -
    \beta
    \sqrt{
      \left(1-\beta^2\right)r_\perp^2
      +
      \mu^2
    }
  }{
    1-\beta^2
  }.
  \label{eq:source_position_parallel_solution}
\end{equation}
  
% ============================================================
\subsection{Evaluation of the Delta Function}
% ============================================================
We have now determined the values of $\vb r'$ for which the argument of the
delta function in the Green function vanishes. At this point one might be
tempted simply to remove the integral and evaluate the integrand at these
values. This, however, is not correct, because the argument of the delta
function is itself a function of the integration variable.

To evaluate the integral we must therefore use the general identity for the
delta function of a vector-valued function,
\begin{equation}
  \delta^3\!\left(\vb F(\vb r')\right)
  =
  \frac{
    \delta^3\!\left(\vb r'-\vb r'_0\right)
  }{
    \left|\det D\vb F(\vb r'_0)\right|
  },
  \label{eq:delta_function_vector}
\end{equation}
where $\vb r'_0$ is the solution of
\[
\vb F(\vb r')=\vb0,
\]
and $D\vb F$ denotes the Jacobian matrix of $\vb F$.

In the present case,
\begin{equation}
  \vb F(\vb r')
  =
  \vb r'
  -
  \vb v
  \left(
    t
    -
    \frac{\abs{\vb r-\vb r'}}{c}
  \right).
  \label{eq:green_function_argument}
\end{equation}

To compute the Jacobian we differentiate each component of $\vb F$ with respect
to the coordinates of $\vb r'$. Since
\[
  \frac{\partial}{\partial r'_i}
  \abs{\vb r-\vb r'}
  =
  -
  \frac{r_i-r'_i}{\abs{\vb r-\vb r'}},
\]
we obtain
\[
  \frac{\partial F_i}{\partial r'_j}
  =
  \delta_{ij}
  -
  \frac{v_i}{c}
  \frac{r_j-r'_j}{\abs{\vb r-\vb r'}}.
\]

Introducing the unit vector
\[
  \vb n
  =
  \frac{\vb r-\vb r'}{\abs{\vb r-\vb r'}},
\]
the Jacobian matrix can be written as
\[
  D\vb F
  =
  I
  -
  \frac{1}{c}\,
  \vb v\otimes\vb n.
\]

Using the identity
\[
\det(I+\vb a\otimes\vb b)=1+\vb b\cdot\vb a,
\]
it follows immediately that
\[
  \det D\vb F
  =
  1
  -
  \frac{\vb v\cdot\vb n}{c}.
\]

The scalar potential therefore becomes
\[
  \phi(\vb r,t)
  =
  q
  \int
  \frac{
    \delta^3\!\left(
      \vb r'
      -
      \vb r'_0
    \right)
  }{
    \left(
      1-\dfrac{\vb v\cdot\vb n}{c}
    \right)
    \abs{\vb r-\vb r'}
  }
  \,d^3r',
\]
where $\vb r'_0$ is the solution of
\[
\vb F(\vb r')=\vb0,
\]
that is, the solution of
\eqref{eq:source_position_condition}. The integration over the delta function
can now be carried out immediately, but the resulting expression still depends
on the unknown source position $\vb r'_0$. Our next task is therefore to
eliminate this dependence.
% ============================================================
\section{The Potentials }
% ============================================================


The integral can now be evaluated immediately, yielding 
\[
  \phi(\vb r,t)
  =
  q
  \frac{
    1
  }{
    \abs{\vb r-\vb r'}
  }
  \,
  \frac{
    1
  }{
    1-\dfrac{\vb v\cdot\vb n}{c}
  },
\]

where it is important to emphasize once again that $\vb r'$ is the solution of
\eqref{eq:source_position_condition}, since this equation is precisely the
condition under which the argument of the function $\vb F$ vanishes.

We now eliminate the dependence on the source position $\vb r'$ so that the
scalar potential is expressed solely in terms of the field point $\vb r$. To
this end, we first compute the scalar product between the particle velocity and
the unit vector $\vb n$. Since the velocity is parallel to the particle
trajectory, only the parallel component contributes, giving

\[
  \vb v\cdot\vb n
  =
  v\,
  \frac{
    r_\parallel-r'_\parallel
  }{
    \sqrt{
      (r_\parallel-r'_\parallel)^2+r_\perp^2
    }
  }.
\]

 the scalar potential becomes
\[
  \phi(\vb r,t)
  =
  q\,
  \frac{
    c
  }{
    c\abs{\vb r-\vb r'}
    -
    v\left(
      r_\parallel-r'_\parallel
    \right)
  }.
\]

From the solution \eqref{eq:source_position_parallel_solution}, we obtain
\[
  r'_\parallel
  -
  v
  \left(
    t
    -
    \frac{\abs{\vb r-\vb r'}}{c}
  \right)
  =
  0.
\]

Multiplying through by $v$ and rearranging gives
\[
  \frac{v^2}{c}
  \abs{\vb r-\vb r'}.
=  v\left(vt-r'_\parallel\right)
\]
 
Finally, using the condition \eqref{eq:source_position_parallel} once more, defining 
\[ \mu = r_\parallel - vt \] 

and after some straightforward algebra, we arrive at the final expression for the
scalar potential, which now depends only on the coordinates of the field point,
\begin{equation}
  \phi(\vb r,t)
  =
  \frac{
    q
  }{
    \sqrt{
      \left(1-\beta^2\right)r_\perp^2
      +
      \mu^2
    }
  }.
  \label{eq:scalar_potential_uniform_motion}
\end{equation}

and the vector potential is, 

\begin{equation}
  {\vb A} = {\vb v} \phi = {\vb v} \frac{
    q 
  }{
    \sqrt{
      \left(1-\beta^2\right)r_\perp^2
      +
      \mu^2
    }
  }.
  \label{eq:vector_potential_uniform_motion}
\end{equation}

\subsection{The Lorentz Gauge}

We have now the expressions for $\phi$ and $\vb A$, but it rests to prove that these solutions fulfill the Lorentz
gauge condition.
Since
\begin{equation}
  \vb A=\frac{\vb v}{c}\phi,
\end{equation}
and $\vb v$ is constant, we have
\begin{equation}
  \nabla\cdot\vb A
  =
  \frac{1}{c}\vb v\cdot\nabla\phi.
\end{equation}

Therefore, the Lorenz gauge condition reduces to
\begin{equation}
  \vb v\cdot\nabla\phi
  =
  -
  \frac{\partial\phi}{\partial t}.
  \label{eq:lorenz_gauge_scalar_potential_condition}
\end{equation}
Since
\begin{equation}
  \phi(\vb r,t)
  =
  q
  \left[
    \left(1-\beta^2\right)r_\perp^2
    +
    \mu^2
  \right]^{-1/2},
\end{equation}
its gradient is
\begin{equation}
  \nabla\phi
  =
  -
  q
  \left[
    \left(1-\beta^2\right)r_\perp^2
    +
    \mu^2
  \right]^{-3/2}
  \left[
    \left(1-\beta^2\right)
    r_\perp\nabla r_\perp
    +
    \mu\nabla\mu
  \right].
\end{equation}

From the decomposition
\begin{equation}
  \vb r
  =
  r_\parallel\vb n
  +
  r_\perp\vb n_\perp,
\end{equation}
it follows immediately that
\begin{equation}
  r_\parallel
  =
  \vb r\cdot\vb n.
\end{equation}
Taking the gradient of both sides gives
\begin{equation}
  \nabla r_\parallel
  =
  \vb n,
\end{equation}
and therefore
\begin{equation}
  \vb v\cdot\nabla r_\parallel
  =
  \vb v\cdot\vb n
  =
  v.
\end{equation}

Since $r_\perp$ is the component of $\vb r$ perpendicular to $\vb n$, we also have
\begin{equation}
  \vb v\cdot\nabla r_\perp
  =
  0.
\end{equation}

From the expression for the gradient,
\[
\nabla\phi
=
-
q
\left[
(1-\beta^2)r_\perp^2+\mu^2
\right]^{-3/2}
\left[
(1-\beta^2)r_\perp\nabla r_\perp
+
\mu\nabla\mu
\right],
\]
and using
\[
\vb v\cdot\nabla r_\perp=0,
\qquad
\vb v\cdot\nabla\mu=v,
\]
we obtain
\[
\vb v\cdot\nabla\phi
=
-
q
\left[
(1-\beta^2)r_\perp^2+\mu^2
\right]^{-3/2}
\mu v.
\]

On the other hand,
\[
\phi
=
q
\left[
(1-\beta^2)r_\perp^2+\mu^2
\right]^{-1/2},
\]
and since
\[
\frac{\partial\mu}{\partial t}
=
-v,
\]
it follows that
\[
\frac{\partial\phi}{\partial t}
=
q
\left[
(1-\beta^2)r_\perp^2+\mu^2
\right]^{-3/2}
\mu v.
\]

Hence,
\[
\vb v\cdot\nabla\phi
=
-
\frac{\partial\phi}{\partial t},
\]
which is precisely the condition
\[
\nabla\cdot\vb A
+
\frac1c
\frac{\partial\phi}{\partial t}
=
0.
\]

% ============================================================
\section{The Electromagnetic Fields}
% ============================================================
 Having obtained both potentials, we can now derive the electric and magnetic
fields.

The electric field is obtained from
\[
  \vb E
  =
  -\nabla\phi
  -
  \frac{1}{c}
  \frac{\partial\vb A}{\partial t}.
\]

Substituting the expressions for $\phi$ and $\vb A$, and after some long algebra,
we obtain
\begin{equation}
  \vb E(\vb r,t)
  =
  q\left(1-\beta^2\right)
  \frac{
    \vb r-vt\,\vb n
  }{
    \left[
      \left(1-\beta^2\right)r_\perp^2+\mu^2
    \right]^{3/2}
  }.
  \label{eq:electric_field_uniform_motion}
\end{equation}

The magnetic field is 
\[
  {\vb B} = \curl{\vb A}
\]
 From
\[
  \vb A=\frac{\vb v}{c}\phi,
\]

Since $\vb v$ is constant,
\[
\curl\vb A
=
\begin{vmatrix}
\vb e_x & \vb e_y & \vb e_z \\
\partial_x & \partial_y & \partial_z \\
v_x\phi & v_y\phi & v_z\phi
\end{vmatrix}
=
  \frac{1}{c}\nabla\phi\times\vb v.
\] 

so 

\[
  {\vb B} =\frac{1}{c}\nabla\phi\times\vb v.
\]

On the hother hand, from the electric field equation above,

\[
  \vb E \times{\vb v}
  =
  -\nabla\phi\times{\vb v}
  -
  \frac{1}{c}
  \frac{\partial\vb A}{\partial t}\times{\vb v} = -c {\vb B}
\]

as $\vb A$ has the same direction of $\vb v$ .

Hence the final formula for the magnetic field is

\begin{equation}
  \vb B(\vb r,t)
  =
   \frac{1}{c} {\vb v} \times {\vb E}
  \label{eq:magnetic_field_uniform_motion}
\end{equation}
 

The electromagnetic field of a point charge moving with constant velocity is
therefore
\begin{align}
  \vb{E}(\vb{r},t)
  &=
  \frac{
    q(1-\beta^2)(\vb{r}-\vb{v}t)
  }{
    \left[
      (r_\parallel-vt)^2
      +
      (1-\beta^2)r_\perp^2
    \right]^{3/2}
  },
  \\
  \vb{B}(\vb{r},t)
  &=
  \frac{1}{c}\vb{v}\cross\vb{E}(\vb{r},t),
  \label{eq:electromagnetics_fields}
\end{align}

% ============================================================
\section{Properties of the Fields} 
% ============================================================

We devote this section to discussing a few important properties of the electric and magnetic fields of a uniformly moving charge.

From \eqref{eq:electromagnetics_fields} it follows immediately that the electric and 
magnetic fields are orthogonal at every point in space and time:

\[
  \vb{E}(\vb{r},t) \cdot  \vb{B}(\vb{r},t) =0
\]

The Poynting vector is given by

\[
\mathbf S
=
\mathbf E\times\mathbf B
=
\frac{1}{c}\,
\mathbf E\times
(\mathbf v\times\mathbf E)
=
\frac{1}{c}
\left[
\mathbf v\,E^2
-
\mathbf E\,(\mathbf E\cdot\mathbf v)
\right].
\]

To simplify the algebra, we introduce the notation

\[
\mathbf E
=
\alpha(\mathbf r-\mathbf vt),
\]

where 
\[
\alpha = \frac{
    q(1-\beta^2) 
  }{
    \left[
      (r_\parallel-vt)^2
      +
      (1-\beta^2)r_\perp^2
    \right]^{3/2}
  },
\]

Substituting this expression into the definition of the Poynting vector yields

\[
\mathbf S
=
\frac{\alpha^2}{c}
\left[
\mathbf v\,(\mathbf r-\mathbf vt)^2
-
(\mathbf r-\mathbf vt)\,
\bigl(\mathbf v\cdot\mathbf r-v^2t\bigr)
\right].
\]

Using the identity

\[
(\mathbf r-\mathbf vt)^2
=
(r_{\parallel}-vt)^2
+
r_{\perp}^{\,2}.
\]

which leads to the final expression for the Poynting vector,

\begin{equation}
  
\mathbf S
=
\frac{\alpha^2 v r_\perp}{c}
\left[
r_\perp\,\mathbf n_\parallel
-
(r_\parallel-vt)\,\mathbf n_\perp
\right].
\label{eq:Poynting_vector}
\end{equation}

Equation \eqref{eq:Poynting_vector} shows that the energy flux is symmetric about the direction of the 
particle velocity, since the direction of the unit vector
 $\vb n_\perp$ is arbitrary within the plane orthogonal to the velocity.
\end{document}
